\begin{solapartat}
\punts{0,2} $A_0 = 0,5\ m$; $f = \dfrac{1}{T} = \dfrac{1}{2}\ Hz$; $\omega = 2\pi f = 2\pi\times\dfrac{1}{2} = \pi\ rad/s$

\punts{0,4} $y(t) = A_0\sin(\omega t + \delta)$

Con que: $y(t = 0) = A_0 \Rightarrow \sin(\omega t + \delta) = 1 \Rightarrow \delta = \pm\dfrac{\pi}{2}$

(també seria vàlida la solució: $y(t) = \pm A_0\cos(\omega t + \delta)$ amb $\delta = 0$)

\punts{0,4} Per tant, l'equació del moviment de la boia és: $y(t) = (0,5\ m)\sin\left(\pi t \pm \dfrac{\pi}{2}\right)$
\end{solapartat}

\begin{solapartat}
\punts{0,2} $v(t) = \dfrac{dy(t)}{dt} = \pm\omega A_0\cos(\omega t + \delta)$

\punts{0,2} La velocitat es màxima quan el $\cos(\omega t + \delta) = 1$\\
$v_{\max} = \pm\omega A_0 = \pm 0,50\pi = \pm 1,6\ m/s$

\punts{0,2} $E_{c,\max} = \dfrac{1}{2}mv^2_{\max}$

% DUBTE: amb v = 1,6 m/s surt 1,92 J; el valor 1,85 J correspon a v = 0,50π = 1,571 m/s sense arrodonir.
\punts{0,4} $E_C = \dfrac{1}{2}1,5\cdot(1,6)^2 = 1,85\ J \approx 1,9\ J$
\end{solapartat}
